% ============================================================
% ANEXO D · ENTORNO DE EJECUCIÓN Y REPRODUCIBILIDAD
% Especificaciones de software, compilación nativa, hardware y semillas.
% Decisiones que lo respaldan: D-007, D-014, D-026.
% ============================================================

\section{Entorno de ejecución y reproducibilidad}\label{anx:entorno}

La reproducibilidad computacional del pipeline depende del control estricto de las versiones
de software, la compilación determinista de extensiones nativas en C y la fijación de semillas
pseudoaleatorias en todas las etapas del procesamiento.

% ------------------------------------------------------------
\subsection{Paquetes y dependencias del sistema}

El desarrollo se articuló en dos entornos independientes de Conda: el entorno \texttt{radiomics},
destinado al preprocesamiento espacial y la extracción de características morfo-texturales
(Módulos 2 y 3), y el entorno \texttt{qml\_cancer}, orientado a la selección estadística,
transformaciones cuánticas variacionales, separabilidad y modelado supervisado (Módulos 4 a 8).
La Tabla~\ref{tab:anexo_d_paquetes} especifica las versiones instaladas de los paquetes fundamentales.

% Generado por Code/tools/tablas_anexos.py desde conda_list_radiomics.txt y conda_list_qml_cancer.txt. No editar a mano.
\begin{table}[H]
\centering
\small
\caption{Versiones de los paquetes de software fundamentales en los entornos de ejecución.}
\label{tab:anexo_d_paquetes}
\begin{adjustbox}{max width=\linewidth}
\begin{tabular}{l p{6.5cm} l l}
\toprule
Paquete de software & Función en el pipeline & Entorno \texttt{radiomics} & Entorno \texttt{qml\_cancer} \\
\midrule
\texttt{python} & Lenguaje base del entorno & \texttt{3.10.20} & \texttt{3.10.19} \\
\texttt{numpy} & Cómputo matricial y vectores & \texttt{2.2.6} & \texttt{2.2.6} \\
\texttt{scipy} & Rutinas científicas y optimización & \texttt{—} & \texttt{1.15.3} \\
\texttt{pandas} & Estructuras de datos tabulares & \texttt{2.3.3} & \texttt{2.3.3} \\
\texttt{scikit-learn} & Modelos clásicos y métricas ML & \texttt{—} & \texttt{1.7.2} \\
\texttt{pyradiomics} & Extracción de características radiómicas & \texttt{0.1.dev1+g8ed579383} & \texttt{3.1.1.dev111+g8ed579383} \\
\texttt{pydicom} & Lectura de mamografías DICOM crudas & \texttt{3.0.2} & \texttt{3.0.1} \\
\texttt{SimpleITK} & Conversión y exportación espacial NRRD & \texttt{2.5.6} & \texttt{2.5.3} \\
\texttt{opencv-python} & Alineación de máscaras y CLAHE & \texttt{4.11.0.86} & \texttt{4.13.0.92} \\
\texttt{matplotlib} & Generación de gráficos y figuras & \texttt{3.10.9} & \texttt{3.10.8} \\
\texttt{torch} & Redes neuronales y optimización MLP & \texttt{—} & \texttt{2.10.0} \\
\texttt{qiskit} & Definición de circuitos cuánticos & \texttt{—} & \texttt{2.3.0} \\
\texttt{qiskit-aer} & Simulación cuántica de alto rendimiento & \texttt{—} & \texttt{0.17.2} \\
\texttt{qiskit-machine-learning} & Kernels y algoritmos cuánticos de aprendizaje & \texttt{—} & \texttt{0.9.0} \\
\bottomrule
\end{tabular}
\end{adjustbox}
\par\smallskip
\parbox{0.95\linewidth}{\footnotesize Versiones exactas de librerías en los dos entornos de ejecución reproducible gestionados con Conda. \texttt{pyradiomics} se encuentra anclado en ambos entornos al commit \texttt{8ed579383b44806651c463d5e691f3b2b57522ab}.}
\end{table}


% ------------------------------------------------------------
\subsection{Procedimiento de compilación de PyRadiomics}

% fuente: D-007
La librería PyRadiomics no admite instalación directa desde los repositorios de PyPI debido a que
la distribución fuente de la versión 3.1.0 declara inconsistencias de versión en sus metadatos y omite
en su manifiesto los archivos de cabecera \texttt{cmatrices.h} y \texttt{cshape.h}, impidiendo la compilación.
Tampoco existe un paquete de PyRadiomics en conda-forge.

Por ello, la biblioteca se compiló directamente a partir del código fuente del repositorio oficial
de AIM-Harvard fijado en el commit \texttt{8ed579383b44806651c463d5e691f3b2b57522ab}. La compilación
se efectuó empleando \texttt{scikit-build-core} y \texttt{setuptools\_scm}, enlazando las extensiones
nativas en C directamente contra el entorno local sin aislamiento de construcción
(\texttt{-\allowbreak-no-\allowbreak build-\allowbreak isolation}) para prevenir actualizaciones accidentales de NumPy que pudiesen comprometer la interoperabilidad con Qiskit.

% ------------------------------------------------------------
\subsection{Especificaciones de hardware y plataforma}

Todas las mediciones temporales, consumos de memoria y ejecuciones de simulación exacta reportados
fueron llevados a cabo en una estación de trabajo homogénea con las siguientes especificaciones:

\begin{itemize}
  \item \textbf{Procesador:} Apple M3 con microarquitectura ARM64 de 8 núcleos (4 núcleos de rendimiento y 4 de eficiencia energética).
  \item \textbf{Memoria principal:} 16\,GB de memoria unificada LPDDR5 (\texttt{hw.memsize}: 17\,179\,869\,184 bytes).
  \item \textbf{Sistema operativo:} macOS Sonoma versión 26.6.2 (Darwin Kernel versión 25G83).
\end{itemize}

% ------------------------------------------------------------
\subsection{Semillas aleatorias y verificación de reproducción determinista}

% fuente: D-014, D-026
Para asegurar la reproducibilidad bit a bit de todos los experimentos se definieron semillas deterministas:
\begin{itemize}
  \item \textbf{Semilla global 42:} empleada para la inicialización pseudoaleatoria del vector de parámetros
    variacionales $\theta$ del \textit{ansatz} (D-014), el sorteo de submuestras estratificadas de 200 lesiones
    para análisis geométrico y el generador de números aleatorios en pruebas de permutación.
  \item \textbf{Semillas de pliegues:} semilla 125 para el subconjunto de masas y semilla 168 para el de
    calcificaciones (D-026), seleccionadas tras una búsqueda exhaustiva entre candidatos para minimizar
    la discrepancia máxima en la prevalencia de malignidad entre pliegues.
\end{itemize}

% fuente: CLAUDE.md, «Re-ejecución completa y revisión (8 al 10 oct)»; H-053, H-054, H-055.
% Corregido el 2026-10-19 (auditoría A2-7): el texto anterior no mencionaba X7 a X9 ni las dos excepciones.
Los Módulos 2 al 8, los experimentos complementarios X1 al X6 y los estudios del simulador B2 y B3 se
re-ejecutaron completos entre el 8 y el 10 de octubre de 2026 sobre un clon limpio del repositorio. Las
predicciones, métricas y tablas se reprodujeron bit a bit, con dos excepciones: los tiempos de ejecución y el
muestreo de B3, que no fija semilla. Los experimentos X7, X8 y X9 se ejecutaron después (10, 11 y 18 de
octubre) con el mismo entorno y las mismas semillas.
